% Figure: the tower crane of the third case study. The assembled machine
% on the left with the two tipping cases marked; the simplified stability
% free-body diagram on the right, reduced to the forces and lever arms
% that enter the overturning balance. Numbers match the worked case study.
\begin{figure}[htbp]
\centering
\begin{tikzpicture}[
  mast/.style={draw=engblack, line width=1.8pt},
  jib/.style={draw=engblack, line width=1.4pt},
  load/.style={draw=engred, -{Stealth}, very thick},
  ballast/.style={draw=engblue, -{Stealth}, very thick},
  dim/.style={draw=enggray, thin, {Stealth}-{Stealth}},
  scale=1.0,
]
% ===== Left: the assembled tower crane =====
\begin{scope}[xshift=0cm]
  % ground
  \draw[draw=engblack, very thick] (-2.6,0) -- (2.8,0);
  \foreach \x in {-2.4,-1.9,...,2.6} {\draw[draw=enggray, thin] (\x,0) -- (\x-0.2,-0.25);}
  % base / foundation block
  \draw[draw=engblack, fill=enggray!18, thick] (-0.9,0) rectangle (0.9,0.5);
  % mast
  \draw[mast] (-0.35,0.5) -- (-0.35,4.0);
  \draw[mast] (0.35,0.5) -- (0.35,4.0);
  \foreach \y in {0.8,1.3,...,3.8} {\draw[draw=enggray, thin] (-0.35,\y) -- (0.35,\y-0.3);}
  % jib (working) to the right, counter-jib to the left
  \draw[jib] (0,4.0) -- (2.4,4.0);
  \draw[jib] (0,4.0) -- (-1.5,4.0);
  % A-frame apex
  \draw[jib] (-0.35,4.0) -- (0,4.6) -- (0.35,4.0);
  \draw[draw=enggray, thin] (0,4.6) -- (2.2,4.05);
  \draw[draw=enggray, thin] (0,4.6) -- (-1.35,4.05);
  % counterweight slab
  \draw[draw=engblack, fill=engblue!18, thick] (-1.5,3.7) rectangle (-1.0,4.0);
  \node[font=\scriptsize, anchor=east] at (-1.5,3.85) {$W_c$};
  % hoist load near jib tip
  \draw[load] (2.2,4.0) -- ++(0,-1.5) node[anchor=north, font=\small, engred] {$W_L$};
  % tipping edges (front and rear base corners)
  \fill[engred] (0.9,0) circle (0.07);
  \node[font=\scriptsize, anchor=north west, engred] at (0.9,0) {$O_f$};
  \fill[engblue] (-0.9,0) circle (0.07);
  \node[font=\scriptsize, anchor=north east, engblue] at (-0.9,0) {$O_b$};
  % radius hints
  \draw[dim] (0,4.25) -- (2.2,4.25);
  \node[font=\scriptsize, anchor=south] at (1.1,4.25) {$R_L$};
  \draw[dim] (0,4.25) -- (-1.25,4.25);
  \node[font=\scriptsize, anchor=south] at (-0.6,4.25) {$R_c$};
  \node[font=\small, anchor=north] at (0.0,-0.45) {(a) the machine};
\end{scope}
% ===== Right: simplified stability FBD about the front tipping edge =====
\begin{scope}[xshift=6.6cm]
  % ground line
  \draw[draw=engblack, very thick] (-2.4,0) -- (2.6,0);
  \foreach \x in {-2.2,-1.7,...,2.4} {\draw[draw=enggray, thin] (\x,0) -- (\x-0.2,-0.25);}
  % schematic crane reduced to a horizontal arm at working height
  \coordinate (Of) at (0.9,0);
  \coordinate (Ob) at (-0.9,0);
  \draw[mast] (0,0) -- (0,3.4);
  \draw[jib] (-1.4,3.4) -- (2.2,3.4);
  % tipping edges
  \fill[engred] (Of) circle (0.07);
  \node[font=\scriptsize, anchor=north, engred] at (Of) {$O_f$};
  \fill[engblue] (Ob) circle (0.07);
  \node[font=\scriptsize, anchor=north, engblue] at (Ob) {$O_b$};
  % load down at jib tip
  \draw[load] (2.2,3.4) -- ++(0,-1.4) node[anchor=west, font=\scriptsize, engred] {$W_L$};
  % counterweight down at counter-jib
  \draw[ballast] (-1.4,3.4) -- ++(0,-1.4) node[anchor=east, font=\scriptsize, engblue] {$W_c$};
  % machine self-weight down at mast centreline
  \draw[draw=engamber, -{Stealth}, very thick] (0,3.4) -- ++(0,-2.0)
    node[anchor=west, font=\scriptsize, engamber] {$W_m$};
  % lever arms from the front edge O_f, drawn along the ground
  \draw[dim] (Of) -- (2.2,0);
  \node[font=\scriptsize, anchor=north] at (1.55,-0.05) {$a_L$};
  \draw[dim] (Of) -- (-1.4,0);
  \node[font=\scriptsize, anchor=north] at (-0.25,-0.05) {$a_c$};
  \node[font=\small, anchor=north] at (0.1,-0.5) {(b) overturning balance about $O_f$};
\end{scope}
\end{tikzpicture}
\caption[A tower crane and its overturning free-body diagram]{The tower
crane of the case study. (a) The machine: a mast carries a working jib to
the right and a shorter counter-jib to the left, with a counterweight
$W_c$ at radius $R_c$ balancing the hoisted load $W_L$ at radius $R_L$.
The base rests on two tipping edges, the front edge $O_f$ (toward the
load) and the back edge $O_b$ (toward the counterweight). (b) The
stability free-body diagram, reduced to the three vertical forces and
their horizontal lever arms about a tipping edge. Forward tipping turns
about $O_f$: the load moment $W_L a_L$ drives the overturn and the
counterweight and machine weight resist it. Backward tipping turns about
$O_b$ with the load removed: the counterweight now drives the overturn
and must itself be resisted. Stability is the statement that the
resisting moment exceeds the driving moment about each edge, with a
margin the standard sets.}
\label{fig:vol03ch01:tower-crane}
\end{figure}
